% \chapter*{KATA PENGANTAR}
% \addcontentsline{toc}{chapter}{KATA PENGANTAR}

% Puji syukur kami panjatkan ke hadirat Tuhan Yang Maha Esa karena atas rahmat dan karunia-Nya, saya dapat menyelesaikan penulisan Laporan Tugas Akhir yang berjudul "Pengembangan Sistem Informasi Berbasis Machine Learning dan Explainable Artificial Intelligence untuk Estimasi Risiko Hipertensi". Laporan ini disusun sebagai salah satu syarat untuk menyelesaikan program Sarjana di Sekolah Teknik Elektro dan Informatika, Institut Teknologi Bandung.
% \begin{flushright}
%     Bandung, 20 Mei 2026\\
    
%     Penulis \\
%     \vspace{1cm}
%     Samuel Franciscus Togar Hasurungan
% \end{flushright}



\chapter*{KATA PENGANTAR}
\addcontentsline{toc}{chapter}{KATA PENGANTAR}

Puji dan syukur senantiasa penulis panjatkan ke hadirat Tuhan Yang Maha Esa atas segala rahmat, kelancaran, serta karunia-Nya, sehingga penyusunan Laporan Tugas Akhir yang berjudul ``Sistem Informasi Berbasis \textit{Machine Learning} dan \textit{Explainable Artificial Intelligence} untuk Estimasi Risiko Hipertensi'' ini dapat diselesaikan dengan baik. Penyusunan laporan ini merupakan salah satu persyaratan akademik guna meraih gelar Sarjana pada Program Studi Sistem dan Teknologi Informasi, Sekolah Teknik Elektro dan Informatika, Institut Teknologi Bandung.

Penelitian dan pengembangan sistem dalam Tugas Akhir ini diharapkan mampu memberikan sumbangsih nyata pada bidang teknologi kesehatan. Secara khusus, sistem ini dirancang untuk menghadirkan pendekatan edukasi naratif yang transparan terkait risiko hipertensi menggunakan \textit{Explainable AI}, memberikan alternatif penyampaian informasi yang lebih terstruktur dibandingkan sistem konvensional. Penulis menyadari bahwa perjalanan dalam merancang dan menyusun penelitian ini penuh dengan dinamika serta proses pembelajaran yang sangat berharga.

Tentu saja, penyelesaian tugas akhir ini tidak lepas dari doa, bimbingan, serta dukungan moral maupun material dari orang-orang di sekitar penulis. Oleh karenanya, penulis ingin menghaturkan rasa terima kasih yang mendalam kepada:

\begin{enumerate}
    \item Keluarga tercinta, khususnya kedua orang tua dan saudari penulis, yang terus memberikan kasih sayang, doa restu, serta motivasi tak terhingga di setiap langkah perjalanan akademik penulis.
    \item Keluarga besar di Medan dan Lubuk Pakam, yang mungkin tanpa penulis sadari juga memberikan dukungan dan doa pada penulis.
    \item Kakak, yang telah menjadi dukungan mental utama penulis selama keberjalanan pengerjaan tugas akhir ini, yang selalu menghibur penulis dengan guyonan abstrak (tidak jelas tapi lucu), selalu mendorong penulis untuk menyelesaikan tugas akhir ini, dan selalu bersedia membantu di saat penulis membutuhkan bantuan teknis.
    \item Mama sebagai partisipan yang telah bersedia meluangkan waktunya untuk meninjau maupun menguji coba sistem yang dikembangkan, atas masukan dan evaluasi yang sangat konstruktif.
    \item Bapak sebagai kepala rumah tangga yang selalu dapat diandalkan dan dengan penuh kesediaan membantu penulis dalam mengurus berbagai keperluan rumah tangga, sehingga penulis dapat lebih memfokuskan waktu dan perhatian pada proses penyelesaian tugas akhir.
    \item Bapak Prof. Dr. Suhono Harso Supangkat, M.Eng. dan Ir. Devi Willieam Anggara, S.T., M.Phil., Ph.D., IPP, selaku Dosen Pembimbing yang telah meluangkan waktu dan tenaganya untuk membimbing penulis dari awal hingga pengerjaan tugas akhir ini dituntaskan.
    \item Seluruh jajaran dosen maupun staf pengajar di lingkungan Program Studi Sistem dan Teknologi Informasi ITB yang telah membekali penulis dengan fondasi ilmu pengetahuan selama masa perkuliahan.
    \item Teman-teman penulis di Purwakarta, yang telah memberikan candaan dan kebersamaan pada penulis selama masa pengerjaan TA.
    \item Playlist Coffee, cafe yang menjadi tempat penulis untuk beristirahat sejenak dan bersosialisasi dengan teman-teman penulis di Purwakarta.
    \item Rekan-rekan seperjuangan dan kolega penulis, atas segala kolaborasi, dukungan, diskusi produktif, serta bantuan teknis yang sangat berarti selama pengerjaan proyek dan masa studi.
    \item Kendrick Lamar dan Kanye West, atas karya-karya musiknya yang telah menemani dan menjadi sumber inspirasi bagi penulis pada saat mengalami kebuntuan ide maupun kekusutan pikiran selama proses pengerjaan tugas akhir.
    % \item Bam Adebayo, yang aksinya bersama \textit{Miami Heat} selalu menjadi sumber hiburan sekaligus inspirasi bagi penulis di saat rehat dari pengerjaan tugas akhir, sebagai bentuk apresiasi atas dedikasi dan performanya yang membuat penulis tetap setia mendukung Miami Heat.
\end{enumerate}

Penulis sangat menyadari bahwa laporan serta sistem yang dibangun ini masih memiliki ruang untuk penyempurnaan dan jauh dari kata sempurna. Maka dari itu, setiap kritik maupun saran yang membangun akan sangat penulis hargai demi perbaikan dan pengembangan penelitian serupa di masa yang akan datang.

Akhir kata, penulis berharap semoga Tugas Akhir ini dapat memberikan manfaat, inspirasi, dan wawasan baru bagi para pembaca, dunia akademik, serta masyarakat luas.

\begin{flushright}
    Bandung, \today\\
    \vspace{1.5cm}
    Penulis, \\
    \vspace{1.5cm}
    Samuel Franciscus Togar Hasurungan
\end{flushright}